\section{引言：从折叠核到“时间作为 uplift”}

本文的出发点是一个最小的、可审计的设定：\textbf{本体是静态约束系统} \(\Clo\)，而非先验给定的演化历史。观察者只能访问部分数据并进行一致性搜索；“时空”与“时间”因此不是输入，而是搜索过程与隐藏自由度组织方式的输出。

\subsection{空间压缩核：\texorpdfstring{$2^m\to F_{m+2}$}{2^m->F_{m+2}}}
在 Zeckendorf 规范中，长度为 \(m\) 的黄金均值语言
\[
\Ym=\{y\in\{0,1\}^m:\ y_i y_{i+1}=0\}
\]
满足 \(|\Ym|=F_{m+2}\)。它可以被视为一个“稳定宏观类型空间”。与之相对，微观读出空间 \(\Xm:=\{0,1\}^m\) 具有 \(2^m\) 个状态。折叠映射 \(\Fold:\Xm\to\Ym\) 把每个微观态压缩到一个宏观稳定态，从而形成空间压缩核。

\subsection{时间作为 uplift：尾部时间纤维}
压缩 \(\Fold\) 一般是多对一，因此若只看 \(\Ym\) 层信息，守恒被打破。本文的关键是：把“丢失的信息”显式放入 uplift 变量，使守恒发生在扩展态上。对 Zeckendorf 截断折叠而言，该 uplift 有一个自然且可计算的实现：\textbf{窗口外的 Zeckendorf 尾部}。本文用尾部移位算子 \(\Tail\) 给出规范时间关系，并把展开过程解释为 \(\Tail^{-1}\) 的分支扩展。

\subsection{能量：并行分支上限与主观单线}
观察者不是全知的：它无法同时追踪所有 uplift 分支。本文用一个最小资源参数 \(\Energy\) 描述这种有限性：\(\Energy\) 是观察者能够维持的并行分支上限（beam width）。当本体约束允许的 uplift 分支数超过 \(\Energy\) 时，观察者必须进行承诺（commit）并丢弃部分分支，从而产生“多世界分支 + 主观单线体验”的资源化解释。

